\section*{Introduction : détermination structurelle et récupération d’informations
}
\addcontentsline{toc}{section}{Introduction : détermination structurelle et récupération d’informations
}
\label{libro:prologo}
L’extrême et moyenne raison holographique pose une question précise : comment une structure peut produire une évaluation numérique, se transformer et conserver l’information nécessaire pour reconnaître et inverser cette transformation. Le problème comprend aussi bien la formation de la valeur que la persistance de ses relations. Une égalité entre quantités permet de vérifier un bilan ; une loi de récupération détermine en outre quel état le réalise, quelles données l’individualisent et par quelle opération il est restitué. Le présent travail développe cette seconde exigence à partir d’une construction arithmétique de chemins et la conduit jusqu’à ses représentations linéaires, incidencielles et physiques.

Le premier résultat qui organise l’argument est la détermination de la constante holographique de couplage arithmético--géométrique. Sa représentation ordonnée est le registre dodécaphasique \(K\) ; son évaluation positionnelle est le rationnel \(\kappa\). La lecture conserve exactement le registre lorsque la base, la longueur et l’origine de phase sont fixées. L’intérêt de l’objet réside également dans ses opérations : ses décalages cycliques forment une base, ses réponses permettent d’identifier des opérateurs et ses images incidentielles transportent les relations de la configuration initiale. Structure, équation et valeur apparaissent ainsi comme des moments d’une construction dont l’information récupérable est déterminée mathématiquement.


Le second résultat fournit la loi générale de cette conservation. Deux transports isométriques avec domaine et codomaine communs produisent deux composantes dont le bilan pondéré restitue exactement le produit scalaire initial. La spécialisation nonadique conserve une norme ternaire de valeur \(73\) et détermine une transformation d'observables avec poids \(72/73\) et \(1/73\). Une politique de prolongement et d'archivage transforme cette identité locale en récupération stable à profondeur arbitraire. La thèse de l'article acquiert ainsi un contenu vérifiable : la valeur évaluée, l'état conservé et la lecture observée sont reliés par des applications explicites et des formules inverses.

\subsection*{Le noyau générateur et ses deux réalisations}
La construction s’établit dans un noyau formel dénommé \emph{holographie modulaire triadique}, ci-après HMT. Son premier objet est \emph{All Possible Paths}, ci-après APP : une structure arithmétique de chemins sur le support toroïdal \((\mathbb Z/9\mathbb Z)^2\), dont le graphe de Cayley décrit les déplacements cardinaux. Sur ce même support, on distingue l’évaluation additive et l’évaluation multiplicative. La conservation du reste avec le quotient maintient l’entier évalué et permet de suivre la provenance de chaque évaluation. La somme accumule ; le produit organise le second feuillet au moyen de sa propre évaluation et de ses règles de transport.


Le second objet, le \emph{TRIT}, apporte la signature orientée \(\{-1,0,+1\}\) et distingue les régimes locaux. L’équation d’Euler étendue exprime, dans ses algèbres quadratiques, les régimes elliptique, parabolique et hyperbolique au moyen d’une loi exponentielle commune. Le troisième objet est le \emph{Topological Phase Kernel}, ci-après TPK. Son action compose sélection, transport et mise à jour de la mémoire. La route, la direction, le feuillet, le résidu, le quotient, la retenue et la frontière sont conservés dans l’état enrichi pertinent ; le retour de phase prolonge son histoire. La connexion et l’holonomie nonadiques appartiennent à cette dynamique.

De cet état proviennent la structure discrète conjointe du continu et ses deux réalisations : la contraction de cylindres compatibles détermine des coordonnées arithmétiques, tandis que l'incidence de frontière conserve des relations combinatoires et géométriques. Les cinq constructions corrélatives du continu demeurent dans ce même développement. Les espaces linéaires, les formes et les observables utilisés ensuite sont des réalisations d'objets déjà engendrés, avec leurs domaines et leurs normalisations. La continuité de l'argument réside dans le fait que chaque opération conserve les relations dont la suivante a besoin.

\subsection*{Complétion de l’unité et raffinement réversible}
La complétion d'une cellule à neuf possibilités par coordonnée ajoute une position de frontière distinguée. En trois coordonnées, la distribution selon le nombre de positions de frontière donne
\[
10^3=729+243+27+1=729+271.
\]
L’intérieur, les strates intermédiaires et le sommet composent une même unité normalisée. En dimension finie arbitraire, le développement de \((9+1)^d\) exprime la même stratification ; la somme des contributions de la coordonnée omise démontre sa compatibilité avec la projection.


Le bloc central \(\bigl(\begin{smallmatrix}7&2\\2&7\end{smallmatrix}\bigr)\) fournit un antécédent arithmétique concret. Ses valeurs propres sont \(9\) et \(5\), et ses concaténations vérifient \(72+27=77+22=99\). La restitution d’une unité et le raffinement suivant forment la chaîne \((72,27)\mapsto(73,27)\mapsto(729,271)\). La première opération incorpore l’unité enregistrée ; la seconde change l’échelle et redistribue une unité au moyen d’une correction de somme nulle. C’est cette suite explicite qui permet de relier centre, complétion et proportion en conservant la provenance de chaque terme.

Le raffinement avec retenue introduit une évolution de cette distribution. La transformation \((x,y)\mapsto(Bx-c,By+c)\) conserve la somme après l'avoir multipliée par \(B\). Sur des étiquettes entières admissibles, avec \(0\leq c<B\), la division euclidienne de la seconde composante récupère d'abord \(c\), puis la paire précédente. La profondeur fixe combien de fois cet inverse doit être appliqué. Lorsque des retenues relevées sont permises, leur quotient supplémentaire fait partie du registre conservé. La réversibilité dépend donc de données précises et possède une opération de récupération explicite.

Cette distinction permet de donner un sens opératoire à la conservation holographique. La lecture normalisée détermine une coordonnée ; le registre conserve les relations qui permettent de reconstruire l’état ou de poursuivre son raffinement. Le développement ultérieur rend cette correspondance linéaire : la bijection entre étiquettes admissibles induit un opérateur unitaire entre leurs bases et permet de transporter amplitudes, corrélations et observables.


\subsection*{La fonction arithmétique et incidencielle de \(K\)}
Le registre \(K\) s’obtient à partir de l’état terminal au moyen de ses visites, traces, dénombrements, différences et charge. L’évaluation
\[
\kappa=\frac{N_K}{1000^{12}-1}
\]
est réversible dans le système de coordonnées fixé : multiplier par le dénominateur restitue l’entier et sa séparation en douze blocs retrouve le registre complet, y compris les zéros initiaux. L’histoire enrichie possède, en outre, les coordonnées de profondeur et de transport qu’utilise chaque prolongement. La récupération de chaque objet est formulée sur les données qui l’individualisent.

Dans la lecture arithmétique, \(\kappa\) intervient dans
\[
\alpha_A=\pi_{\rm HMT}+e_{\rm HMT}-\varphi_{\rm HMT}-4-\kappa.
\]
Les trois coordonnées régionales et le registre procèdent du même noyau générateur et agissent selon leur ordre constructif. Le système de coordonnées analytique complet de \(\alpha\) représente cette même précoordonnée au moyen d’une racine simple et d’une collection compatible d’intervalles. Dans la lecture incidentielle, le registre transporte ses relations de phase vers la configuration exceptionnelle. L’unité de l’objet se manifeste dans la composition effective de ses lecteurs, de ses inverses et de ses actions.

L’orbite cyclique de \(K\) ajoute un résultat d’identification. Ses douze vecteurs forment une base de la réalisation linéaire à douze composantes. Les réponses à cette base déterminent tout opérateur linéaire de la classe considérée ; pour un opérateur qui commute avec le décalage, une réponse vectorielle complète détermine les autres. La plus petite valeur propre du Gram, \(589609\), fournit une borne exacte de stabilité. Deux observations complémentaires de mémoire retrouvent également la partie centrée du registre ; la charge uniforme complète sa reconstruction. La constante fonctionne ainsi comme un objet de couplage doté d’une capacité quantifiée d’évaluation, de transport et d’identification.

\subsection*{Le bilan bilatéral et la transformation des observables
}
L’entier \(73\) apparaît comme norme de deux éléments conjugués et comme indice d’une inclusion réticulaire de rang deux. L’identité \(10\cdot73=729+1\) relie cette norme ternaire à la complétion décimale. Le théorème bilatéral étend le bilan à deux isométries \(B,C\) entre espaces ayant un domaine et un codomaine communs. Pour \(a>0\), les composantes \(aB-C\) et \((a+1)B+C\), pondérées par \(a+1\) et \(a\), conservent le produit scalaire après la normalisation indiquée. Le développement et l’annulation des termes croisés prouvent l’identité en dimension arbitraire. La spécialisation \(a=8\) maintient les opérateurs nonadiques et leur information arithmétique.


La conservation de l’état complet permet de calculer l’effet d’une observation particulière. Si le transport relatif \(R\) est unitaire et \(G\) est l’observable exprimé dans le système de coordonnées commun, la lecture diagonale transforme
\[
\mathcal T(G)=\frac{72}{73}G+\frac1{73}R^*GR.
\]
Un observable reste invariant exactement lorsqu’il commute avec \(R\). L’observable transporté sur l’espace complet conserve toute la lecture initiale. Le résultat précise, dans une même construction, l’information qui persiste et le changement introduit lorsqu’on l’observe à travers un canal déterminé.

\subsection*{Récupération à profondeur arbitraire}
À chaque étape, on peut prolonger une composante et archiver son complément. Le bilan fini additionne les carrés des normes des mémoires et de l'état terminal. Une politique explicite impose la borne géométrique \((729/1241)^N\) sur le carré de la norme de l'opérateur terminal nonadique. À la limite, toute la norme est contenue dans l'archive infinie ; cette archive est une isométrie et son adjoint reconstruit l'état avec une norme d'opérateur égale à un.

Les premiers blocs archivés admettent également un inverse exact corrigé par leur opérateur de Gram, avec une borne de stabilité uniforme pour les profondeurs positives. L’archive infinie, l’inverse fini et l’approximation obtenue en tronquant l’adjoint possèdent des formules et des estimations propres. Dans les dynamiques générales où un terminal subsiste, sa contribution est conservée avec la mémoire. La politique démontrée identifie le cas où ce terminal tend vers zéro.

Le relèvement unitaire du raffinement se compose avec le bilan bilatéral. L’exemple décimal qui conduit \((73,27,1)\) à \((729,271)\) maintient explicite la provenance de la retenue ; sa comparaison avec une retenue voisine détermine le changement exact d’une lecture. La conservation évolutive de l’unité s’exprime par une séquence de transformations récupérables.


\begin{figure}[htbp]
\centering
\begin{tikzpicture}[>=Stealth,every node/.style={font=\small,align=center},
block/.style={draw=ink,rounded corners=2pt,inner sep=6pt,text width=35mm},
wide/.style={draw=ink,inner sep=6pt,text width=47mm}]
\node[block] (d) {État arithmétique\\admissible\\\((x,y,c)\)};
\node[block,right=14mm of d] (e) {État raffiné\\\((10x-c,10y+c)\)};
\node[block,right=14mm of e] (j) {Bases d’étiquettes\\\(J\) unitaire};
\draw[<->] (d)--node[above,font=\scriptsize]{retenue} (e);
\draw[->] (e)--(j);
\node[wide,below=17mm of e] (w) {Transport bilatéral\\\(W^*W=I\)};
\draw[->] (j.south)|-(w.east);
\node[block,below left=17mm and 4mm of w] (a) {Lecture conjointe\\\(\frac{72G+R^*GR}{73}\)};
\node[block,below right=17mm and 4mm of w] (b) {Observable transporté\\\(WGW^*\)};
\draw[->] (w)--(a);\draw[->] (w)--(b);
\node[font=\small,below=8mm of a] {Invariant si \(GR=RG\)};
\node[font=\small,below=8mm of b] {Récupération au moyen de \(W^*\)};
\end{tikzpicture}
\caption{Composition d’opérations démontrées. Le raffinement conserve la somme normalisée et permet de récupérer l’étiquette. Son relèvement agit sur les amplitudes des états et conserve la norme. La lecture d’un observable et son transport complet conservent des fonctions différentes, spécifiées par les deux formules inférieures.}
\label{libro:fig:composicion}
\end{figure}


\subsection*{Formes, action et observation quantique}
Le transport du produit scalaire se prolonge à chaque degré extérieur admissible. Aires et volumes conservent également les contributions mixtes de lecture et de mémoire. L'incidence exceptionnelle transporte ces identités sur ses images spécifiées ; les charges variationnelles s'obtiennent à partir d'une action déclarée et de ses symétries. Les réalisations électromagnétiques et gravitationnelles maintiennent les échelles et les opérateurs qui interviennent dans leurs preuves.

La réalisation gaussienne montre pourquoi conserver l’état et connaître une observation réduite sont des questions complémentaires. Dans le protocole déclaré de deux entrées gaussiennes pures, centrées et identiques et d’un nombre fini de modes, le transport symplectique conserve la structure globale, tandis que sa réduction acquiert un spectre déterminé par la composante antilinéaire de la mémoire. La boucle elliptique--parabolique étudiée produit une pureté \(9/11\), un spectre \((9/10)10^{-j}\) et une entropie adimensionnelle \((10/9)\log10-\log9\), avec toutes les occupations de l’espace de Fock conservées dans le calcul.


L’argument s’achève donc sur une loi de transformation plus riche qu’une égalité scalaire : la valeur est déterminée, la structure qui intervient dans sa lecture est conservée et les conditions de sa récupération sont calculées. Les formules d’inversion, les bornes de stabilité et les réalisations spécifiées expriment conjointement le sens mathématique de l’extrême et moyenne raison holographique.

% BEGIN SINTESIS_ADITIVA_20260922
\par
La section~\ref{delta22:xi:holonomia} explicite la contribution de l'archive à la composition bilatérale, construit le plan de quart de tour du calendrier dodécaphasique et démontre la restitution orientée de l'action au moyen d'un commutateur relatif. L'amplification compatible conserve l'opérateur ; le détecteur distingue la dilatation de l'état du gain énergétique et l'archive ordonnée du décompte des opérations.
% END SINTESIS_ADITIVA_20260922
